\documentclass[../geometrie.tex]{subfiles}

\begin{document}
	\subsection{Espace vectoriel quotient}{
		A compléter plus tard.
	}

	\subsection{Dual topologique}{
		\noindent Soit $V$ un $\mathbb{R}$-espace vectoriel normé. \\
		\noindent Soit $\mathcal{B}$ une base de $V$.

		\vspace{1em}

		On s'intéresse ici au dual topologique de $V'$ qui s'avèrera utile par la suite, en particulier lors de preuve de liberté de certaines familles.

		\mathblock{definition}{Dual topologique}{
			On appelle \textit{dual topologique} de $V$ l'ensemble noté $V'$ définie par
			$$
				V' := \mathcal{L}^{c}(V, \mathbb{R})
			$$
			qui correspond à l'ensemble des formes linéaires continues de $V$.
		}

		On s'intéresse à présent à certaines formes linéaires particulière de $V'$, à savoir celle \flqq{} engendrées \frqq{} à partir d'éléments d'une base de $V$.

		\mathblock{definition}{Application $b^{*}$}{
			Soit $b_{0} \in \mathcal{B}$. \bigskip \\
			On notera $b_{0}^{*}$ l'application définie par
			$$
				\forall b \in \mathcal{B}, b_{0}^{*}(b) := \delta_{b_{0}, b} =
				\left\{
					\begin{array}{lcl}
						1 & \text{si} & b = b_{0} \\
						0 & \text{si} & b \neq b_{0}
					\end{array}
				\right.
			$$
		}

		Pour préciser le comportement de $b_{0}^{*}$, on aura donc pour $v \in V$ l'expression suivante de $b_{0}^{*}(v)$. \\
		Soit l'unique famille presque nulle $(\lambda_{b})_{b \in \mathcal{B}} \in \mathbb{R}^{\mathcal{B}}$ telle que
		$$
			v = \sum_{b \in \mathcal{B}} \lambda_{b} \times b
		$$
		Alors
		$$
			b_{0}^{*}(v) = \lambda_{b_{0}}
		$$
	}
\end{document}
